% !TEX root = userguide.tex

\chapter{Stochastic fluid models}
\pgst{tfemchA}
\label{chap:stochasticmodels}

\def\chaptertitle{Stochastic fluid models}
\def\headdate{6th February 2009}

In this Appendix some information is provided on the stochastic fluid
models available in the add-on `bcf'.

\section{Notation}

The Cauchy stress tensor is denoted by $\ten\sigma$. We assume that
the fluid is incompressible.  Hence, we split $\ten\sigma$
as follows
\[
   \ten\sigma = -p\ten 1 +\ten\tau_\text{ex},
\]
where $p$ is the pressure and $\ten\tau_\text{ex}$ is the extra-stress tensor,
which vanishes in equilibrium.
The tensor $\ten\tau_\text{ex}$ is split into a solvent part and a polymer part
\[
   \ten\tau_\text{ex} = 2\eta_s\ten D + \ten\tau, 
\]
where $\eta_s$ is the solvent viscosity and $\ten D$ is Euler's
rate-of-deformation tensor: 
\[
         \ten D=\frac12(\ten L + \ten L^T)\quad\text{with}\quad 
         \ten L^T= \nabla\vek u.
\]

\section{Dumbbell model}

This model is based on elastic dumbbells. The Langevin equation is
given by
\begin{equation}
  \mathrm{d}\vek{Q}=(\ten L\cdot\vek Q - \frac{1}{2\lambda_H}\vek F_c)
    \mathrm{d}t+\frac{1}{\sqrt{\lambda_H}}\mathrm{d}\vek W,
\label{eq:Langevin}
\end{equation}
where $\vek Q$ is the dumbbell connector vector scaled with
$\sqrt{kT/H}$, $\lambda_H=\zeta/4H$ is a time constant, $\mathrm{d}\vek W$
is a Wiener increment and $\vek F_c$ is the connector force given by
\begin{align*}
     \text{Hookean dumbbells:\qquad}\vek F_c&=\vek Q,\\
     \text{FENE-P:\qquad}\vek F_c&=\frac{\vek Q}{1-\langle Q^2\rangle/b}.\\
     \text{FENE:\qquad}\vek F_c&=\frac{\vek Q}{1-Q^2/b}.
\end{align*}
The expression for the stress tensor is
\begin{equation}
 \ten \tau_p=G(\langle\vek Q\vek F_c\rangle-\ten 1 ),
 \label{eq:sttress}
\end{equation}
with $G=nkT$.

In equilibrium the structure tensor $\ten\alpha=\langle\vek Q\vek Q\rangle$
is given by
\begin{align*}
     \text{Hookean dumbbells:\qquad}\ten\alpha_{\text{eq}}&=\ten 1,\\
     \text{FENE-P:\qquad}\ten\alpha_{\text{eq}}&=\frac{b}{b+3}\ten 1,\\
     \text{FENE:\qquad}\ten\alpha_{\text{eq}}&=\frac{b}{b+5}\ten 1.
\end{align*}


\section{Filling of the material parameters}

The properties of the model are stored in a structure of type
\texttt{stmodel\_t}. The structure must first be constructed by calling the
routine \texttt{create\_stochastic\_model}. Then the material parameters must
be filled in the components \texttt{modulus}, \texttt{lambda}, \texttt{nonlin}
of the structure. In Table~\ref{tab:matparbcf} it is shown which parameters need
to be filled for each model
\begin{table}[!hbtp]
\begin{center}
   
  \caption{Material parameters}\label{tab:matparbcf}\medskip
  \begin{tabular}{l|ccc}
    \texttt{model} & \texttt{modulus} & \text{lambda} & \texttt{nonlin}\\\hline
          1 (Hookean dumbbell) & $G=nkT$ & $\lambda_H$ & - \\
          2 (FENE-P)  &  $G=nkT$ & $\lambda_H$ & $b$ \\
          3 (FENE)    &  $G=nkT$ & $\lambda_H$ & $b$ 
  \end{tabular}

\end{center}
\end{table}

\section{Implementation}

All the models are implemented with three component equations
($Q_x$, $Q_y$ and $Q_z$), except for the Hookean dumbbell in planar flows,
which use only two ($Q_x$ and $Q_y$). The number of stress and conformation
tensor components is four in 2D/axisymmetrical flows and six in 3D,
except the for Hookean dumbbell in planar flows,
which has only three.

The Langevin equation is solved using Brownian dynamics. The numerical
parameters for this method should be stored in a structure of type
\texttt{stnumpar\_t}, with components \texttt{nfield} (number of
fields/realizations), \texttt{timeint} (the type of integration)
and \texttt{tstep} (the time step).
